%% Teste de integração com bibliografia
%% Este arquivo testa os diferentes sistemas de bibliografia da classe UnB-PPGT

\documentclass[mestrado,natbib]{UnB-PPGT}

%% Metadados mínimos para teste
\titulo{Teste de Integração Bibliográfica}
\autor{Teste}{Usuario}
\orientador{Prof. Dr. Orientador Teste}{Universidade de Brasília}
\diamesano{1}{1}{2024}
\membrobancafunc{Prof. Dr. Orientador Teste}{Universidade de Brasília}{Orientador}
\membrobancafunc{Prof. Dr. Membro 1}{Universidade de Brasília}{Examinador}
\membrobancafunc{Prof. Dr. Membro 2}{Universidade Federal de Goiás}{Examinador Externo}

\begin{document}

\pretextual

\begin{resumo}
Este é um teste da integração bibliográfica da classe UnB-PPGT. O sistema deve suportar diferentes formas de citação e formatação de referências.
\end{resumo}

\textual

\chapter{Teste de Citações}

Este capítulo testa diferentes tipos de citação disponíveis na classe UnB-PPGT.

\section{Citações Básicas}

Exemplo de citação no texto: \citeonline{vasconcellos2012} apresenta uma análise da mobilidade urbana.

Exemplo de citação entre parênteses: A mobilidade urbana é um tema complexo \cite{vasconcellos2012}.

\section{Citações com Página}

Citação com página específica: \citepag{p. 45}{silva2019}.

Citação no texto com página: \citeonlinepag{p. 123}{oliveira2021} discute sustentabilidade.

\section{Citações Múltiplas}

Vários autores tratam do tema \cite{vasconcellos2012,silva2019,oliveira2021}.

\section{Citações Especiais}

Citação apud: \citeapud{costa2018}{silva2019}.

Citação de legislação: \citelaw{brasil2012}.

Citação de dados: \citedata{ibge2022}.

\postextual

\bibliography{bibliografia}

\end{document}